% TOP_PROBLEM: {"id": 1200044, "title": "Some Questions — Question 44", "category_id": 17, "category": {"id": 17, "name": "group_theory", "display_name": "Group Theory", "description": "Problems about groups, group actions, representations, and related algebraic structures.", "slug": "group-theory", "order_index": 17, "created_at": "2026-07-31T15:26:25.671Z"}, "status": "open", "problem_number": "AMR-011-0044", "set_id": 13}
% FIRST_POSTED: 2026-10-01
% ENTRY_KIND: statement_only
% UnsolvedMath ID 1200044; source catalogue code AMR-011-0044
% !TeX program = lualatex
% Definition and sources only; this file is not an LLM solution attempt.
\documentclass[11pt,a4paper]{article}
\usepackage[margin=25mm]{geometry}
\usepackage{fontspec}
\setmainfont{lmroman10-regular.otf}[BoldFont=lmroman10-bold.otf,
 ItalicFont=lmroman10-italic.otf,BoldItalicFont=lmroman10-bolditalic.otf]
\usepackage{amsmath,amssymb,mathtools}
\usepackage{unicode-math}
\setmathfont{latinmodern-math.otf}
\AtBeginDocument{\let\setminus\smallsetminus}
\usepackage{xurl,hyperref}
\hypersetup{colorlinks=true,urlcolor=blue}
\setcounter{secnumdepth}{0}
\setlength{\parindent}{0pt}
\setlength{\parskip}{5pt}
\setlength{\emergencystretch}{3em}
\begin{document}
\section*{Some Questions — Question 44}\label{1200044}
{\small UnsolvedMath ID 1200044 · AMR-011-0044 · Group Theory · AMR Open Problem Lists}
\subsection{Definitions and mathematical statement}
Let a countable group $\Gamma$ act essentially freely by probability-measure-preserving Borel transformations on a standard atomless probability space $(X,\mu)$. Its orbit relation $R_X$ relates $x,y$ when $y=\gamma x$ for some $\gamma\in\Gamma$. A graphing is a countable family of measure-preserving partial bijections between measurable subsets of $X$ whose generated equivalence relation is $R_X$ modulo null sets. Its cost is the sum of the measures of the domains, and the cost of the action is the infimum over all such graphings.

Give $X\times X$ the product probability measure and the diagonal action
\[
 \gamma\cdot(x,y)=(\gamma x,\gamma y).
\]
The same group element acts in both coordinates; this is not an action of $\Gamma\times\Gamma$. It is again essentially free.

Must
\[
 \operatorname{cost}(\Gamma\curvearrowright X)
 =\operatorname{cost}(\Gamma\curvearrowright X\times X)
\]
hold for every such group and action? The equality includes extended-real values if an action has infinite cost. No ergodicity assumption is imposed on the original action or its diagonal square.

In particular, the orbit relation of the diagonal action is not the full product relation obtained by moving the coordinates independently. The question compares the infima over all graphings of these two specific relations, not the costs of particular lifted generators. The projection onto either factor is equivariant, but the existence of that map does not by definition identify the two cost invariants.
\subsection{Short English statement}
Does passing from a free probability-measure-preserving action to its diagonal square preserve cost?
\subsection{Sources}
\begin{itemize}
\item[S1] ULAM.AI and the UnsolvedMath Contributors, supplied problems.json, record 1200044 (AMR-011-0044), AMR Open Problem Lists. Catalogue identity and source transcription; CC BY 4.0.
\url{https://huggingface.co/datasets/ulamai/UnsolvedMath}.
\item[S2] Abert - Some questions (pdf) (2010), Question 44. Original source indexed by AMR-011-0044.
\url{https://www.renyi.hu/~abert/questions.pdf}.
\end{itemize}
\end{document}
